%======== BEGIN inserted by add_gemeinsam.rb ======================
\documentclass{beamer}
\usepackage[german]{babel}
\usepackage[T1]{fontenc}
%\usepackage[latin1]{inputenc} 
\usepackage[utf8x]{inputenc}

\usepackage{verbatim} 
\usepackage{listings} 

% \lstset{numbers=left, numberstyle=\tiny, numbersep=5pt} 
\logo{\includegraphics[width=0.7cm]{../images/ruby.png}}

%\usetheme{Warsaw}
%\usetheme{Hannover}
%\usetheme{PaloAlto}
%\usetheme{JuanLesPins}
%\usetheme{Antibes}
%\usetheme{Shingara}
%\usetheme{Berlin}
%\usetheme{Oxygen}
\usetheme{Singapore}

\usecolortheme{beaver}
% albatross | beaver | beetle |
% crane | default | dolphin |
% dove | fly | lily | orchid |
% rose |seagull | seahorse |
% sidebartab | structure |
% whale | wolverine


\title[Ruby]{Programmiersprache Ruby}
\author{Sven Suska, Thomas Preymesser}
%\institute{}
\date{2009-Juli-2}


\usepackage{tikz}
\usetikzlibrary{arrows,shapes,decorations}


\tikzstyle{klasse} = [rectangle, draw, thin, fill=blue!20 ]
\tikzstyle{objekt} = [ellipse, draw, thin, fill=green!20, minimum height=2.5em]

\tikzset{%
    class/.style={green,>=triangle 45}
    }
\tikzset{%
    eigenc/.style={very thin,green,>=triangle 45}
    }
\tikzset{%
    superc/.style={thick,red,>=open triangle 45}
    }
\tikzset{%
    method1/.style={thin,blue,bend left=45,to reversed->}
    }
\tikzset{%
    method1r/.style={thin,blue,bend right=45,to reversed->}
    }
\tikzset{%
    methods/.style={thin,blue,bend left=45}
    }

\newcommand{\klassebox}[2]{\begin{tabular}{l}{\it #1} \\ \hline \tt{#2}\end{tabular}}

\newcommand{\bls}[1]{{\bf\lstinline{#1}}}

\begin{document}
\lstset{language=Ruby}
\lstset{basicstyle=\small,numbers=none, numberstyle=\tiny, numbersep=5pt}
%======== END inserted by add_gemeinsam.rb ======================
\begin{frame}
  \frametitle{Threads}
  \lstinputlisting[frame=single,caption={Kuchen backen}]{code/kuchenbacken.rb}
\end{frame}

\begin{frame}
  \frametitle{Threads}
  \lstinputlisting[frame=single,caption={Kuchen backen Ausgabe}]{code/kuchenbacken_ergebnis}
\end{frame}

\begin{frame}
  \frametitle{Threads}
  \begin{itemize}
    \item Jeder Thread l"auft parallel zum Hauptprogramm
    \item Ruby Threads sind Threads innerhalb des Interpreters
    \begin{itemize}
      \item + portabel auf allen Systemen
      \item - keine Mehrprozessor-Systeme unterst"utzt
      \item - System-Calls blockieren andere Threads im Interpreter
    \end{itemize}
  \end{itemize}
\end{frame}



\begin{frame}
  \frametitle{Thread Variablen}
  Variablen in Threads
  \begin{itemize}
    
    \item externe Variablen
    \item Thread-lokale Variablen
    \item Thread Austausch Variablen \\
          \texttt{Thread.current['zeit'] = 15}
  \end{itemize}
\end{frame}

\begin{frame}
  \frametitle{Exceptions in Threads}
  Exceptions in Threads
  \begin{itemize}
    \item \texttt{Thread.abort\_on\_exception = false} (default)\\
       nicht behandelte Exception beendet diesen Thread
    \item \texttt{Thread.abort\_on\_exception = true}\\
       nicht behandelte Exception beendet alle Threads
  \end{itemize}
\end{frame}


\begin{frame}
  \frametitle{Threads}
  Kontrolle des Thread-Schedulers
  \begin{itemize}
    \item \texttt{Thread.stop} 
    \item \texttt{Thread\#run}
    \item \texttt{Thread.pass}
  \end{itemize}
\end{frame}

\begin{frame}
  \frametitle{Threads}
  \lstinputlisting[frame=single,caption={Kuchen backen}]{code/kuchenbacken2.rb}
\end{frame}

\begin{frame}
  \frametitle{Threads}
  Monitors
  \lstinputlisting[frame=single,caption={counter.rb}]{code/counter.rb}
\end{frame}

\begin{frame}
  \frametitle{Threads}
  Konkurrierende Zugriffe auf Resource
  \begin{itemize}
    \item Problem: \texttt{@counter += 1}\\
    \pause
    \item L"osung: Monitor
  \end{itemize}
\end{frame}

\begin{frame}
  \frametitle{Threads}
  Monitors
  \lstinputlisting[frame=single,caption={counter2.rb}]{code/counter2.rb}
\end{frame}

\begin{frame}
  \frametitle{Threads}
  Monitor nicht als Unterklasse von Monitor
  \lstinputlisting[frame=single,caption={counter3.rb}]{code/counter3.rb}
\end{frame}

\begin{frame}
  \frametitle{Threads}
  ThreadGroup ???
\end{frame}

%======== BEGIN inserted by add_gemeinsam.rb ======================
\end{document}
%======== END inserted by add_gemeinsam.rb ======================
